\documentclass[11pt]{article}

\usepackage[a4paper,margin=1.1in]{geometry}
\usepackage{amsmath,amssymb,amsthm,mathtools}
\usepackage{tikz}
\usepackage{tikz-cd}
\usepackage{microtype}
\usepackage{enumitem}
\usepackage{booktabs}
\usepackage[hidelinks]{hyperref}
\usepackage{xcolor}
\usepackage{listings}

\definecolor{accent}{RGB}{0,82,147}
\definecolor{commentgray}{RGB}{110,110,110}
\hypersetup{colorlinks=true,linkcolor=accent,citecolor=accent,urlcolor=accent}

\lstdefinestyle{sketch}{
  basicstyle=\ttfamily\footnotesize,
  commentstyle=\color{commentgray}\itshape,
  keywordstyle=\color{accent}\bfseries,
  morecomment=[l]{\#},
  morekeywords={bridge,shared,entity,op,cone,legs,into,annotations,lost,partial,provenance,attested,last\_verified},
  columns=fullflexible,
  keepspaces=true,
  frame=single,
  framesep=6pt,
  rulecolor=\color{commentgray},
  xleftmargin=2pt,xrightmargin=2pt,
  aboveskip=1em,belowskip=0.5em,
}

\theoremstyle{plain}
\newtheorem{theorem}{Theorem}
\newtheorem{proposition}{Proposition}
\newtheorem{lemma}{Lemma}
\newtheorem{corollary}{Corollary}
\theoremstyle{definition}
\newtheorem{definition}{Definition}
\newtheorem{example}{Example}
\newtheorem{remark}{Remark}

\newcommand{\Set}{\mathbf{Set}}
\newcommand{\Cat}{\mathbf{Cat}}
\newcommand{\Mod}{\mathrm{Mod}}
\newcommand{\Sketch}{\mathbf{Sk}}
\newcommand{\Lex}{\mathrm{Lex}}
\newcommand{\op}{^{\mathrm{op}}}
\newcommand{\Trade}{\mathsf{Trade}}
\newcommand{\Inst}{\mathsf{Instrument}}
\newcommand{\FinProd}{\mathsf{FinancialProduct}}
\newcommand{\Client}{\mathsf{Client}}
\newcommand{\Wallet}{\mathsf{Wallet}}
\newcommand{\Cpty}{\mathsf{Counterparty}}
\newcommand{\Cust}{\mathsf{Customer}}
\newcommand{\restrict}[1]{\!\!\upharpoonright_{#1}}

\title{\textbf{The Enterprise Model Is a Colimit}\\[0.3em]
\large Compositional Integration of AI-Readable Domain Ontologies}

\author{Iman Poernomo\\[0.25em]\normalsize Bullish\\[0.1em]\normalsize\texttt{iman.poernomo@gmail.com}}
\date{June 2026}

\begin{document}
\maketitle

\begin{abstract}
The ``unified enterprise data model'' (EDM), in which one global schema renders every domain a sub-domain of a consistent whole, is the oldest unkept promise in enterprise data. Its failure is usually attributed to a moving target: the enterprise is said to change faster than any model can track. The core business domains of a trading firm change slowly, on the timescale of years, and the EDM still never ships, so the moving-target account is inadequate and the failure is structural. The EDM is a colimit of a diagram of domain theories glued along what they share, and the historical error was to author that colimit as a single complete artifact instead of assembling the diagram incrementally, where each stage is already useful. We model each domain as a finite-limit sketch, a bridge between two domains as a span through a shared sub-theory, two-domain integration as a pushout, and multi-domain integration as a colimit. We prove three results. Consistency of an integration is equivalent to commutativity of its bridge diagram, and a violation localizes to one witnessing concept and one pair of disagreeing bridges, which makes ``is our integration correct?'' a decidable check that returns a counterexample. The model functor sends the theory-level colimit to the corresponding limit of instance categories, so an instance of the integrated model is a compatible tuple of domain instances agreeing on shared sub-theories; a governance rule we state---\emph{if two domains genuinely share a concept their rules agree, otherwise split the concept}---is the side condition under which this holds strictly. And every sub-diagram embeds faithfully into the whole, so partial assemblies stay correct and are never invalidated by later gluing, and global completeness is never required. The mathematics is the specification-gluing of \cite{cpw2000,pwc2005}, here transported from algebraic specifications to AI-readable domain ontologies. The contribution is the transport, the diagnosis of the EDM and the OWL/RDF program, and the observation that LLM-legible prose-anchored sketches make the objects cheap to author where description logic made them prohibitive.
\end{abstract}

\medskip
\noindent\textbf{MSC 2020.} 18C30 (sketches and generalizations); 18A30 (limits and colimits); 18C10 (theories, structure, and semantics); 68P15 (database theory).

\noindent\textbf{ACM CCS.} Information systems $\rightarrow$ Data management systems; Theory of computation $\rightarrow$ Logic (categorical semantics).

\noindent\textbf{Keywords.} colimit; pushout; finite-limit sketch; ontology integration; enterprise data model; institutions; semantic layer.

\section{Introduction: the model that never ships}

Every large enterprise has attempted, at least once, to build a single coherent model of its data: one vocabulary, one set of canonical entities, every domain a tidy sub-domain of a consistent whole. The aspiration has many names (the enterprise data model, the logical data model, the corporate information factory, the canonical ontology) and a uniform outcome. Teams of senior, expensive, capable people work for months or years; artifacts are produced and presented; and the model is either never completed or is completed for an organization that has since moved on. In financial services the pattern is visible at the scale of decades and industry consortia: standardized vocabularies for which adoption beyond a center of excellence remains, in the field's own assessments, persistently out of reach.

The received explanation is that the enterprise is a moving target: domains evolve, acquisitions arrive, regulations redefine entities, and no hand-maintained model can keep pace. This account locates the failure outside the method, which is part of its appeal, and it does not survive contact with the timescales involved. The core domains of a trading business (what a trade is, what a client is, what distinguishes an instrument from a financial product, what a settlement is) are stable, changing over years rather than weeks. A standardized derivatives vocabulary can remain substantially fixed, and substantially unadopted, for two decades. Slow-moving targets are the ones a hand-maintained model should have been able to hit, and the EDM failed to hit them, which places the failure in the method rather than in the rate of change.

\paragraph{Thesis.} The unified enterprise model is a \emph{colimit}: the universal object obtained by gluing a diagram of domain theories along the concepts they share. A colimit of a slowly-changing diagram does not fail to exist because the diagram drifts. It fails to \emph{ship} when one insists on presenting the colimit directly, as a single globally authored internally consistent monolith, instead of presenting the \emph{diagram} and letting the colimit be whatever the parts and their gluings determine. Direct presentation has two structural costs. It admits no partial value, since a half-finished global schema is a model of nothing until it is finished. And it places correctness in a global consistency check on the merged monolith rather than a local check on whether the bridges agree, so a defect implicates the whole model and cheap termination is lost. Presenting the colimit as a diagram avoids both: every sub-assembly is immediately usable, and correctness reduces to the commutativity of finitely many bridge paths.

This reframing also accounts for the most rigorous failed attempt at the problem. The Semantic Web program (RDF, OWL, linked data) adopted the open-world assumption, under which meaning composes and no global completeness is demanded, which is the right starting point. It then encoded that assumption as description logic over a merged global graph, which reinstated the monolith by making correctness a property of the whole, and it chose a serialization that humans, and later language models, could neither author nor read at scale. The categorical framing retains the open-world locality and drops the description-logic monolith. The objects to be glued are then no longer first-order theories in awkward markup but domain ontologies written as structured, prose-anchored sketches that a capable language model reads directly. The machinery that made the open-world assumption prohibitively expensive is the machinery this paper removes.

\paragraph{Provenance of the construction.} We claim no new mathematics. Gluing theories along shared sub-theories by pushout, and composing many along a diagram by colimit, is the specification-union construction the author developed for algebraic specifications and their proofs \cite{cpw2000,pwc2005}, within the tradition of institutions \cite{goguen-burstall} and structured algebraic specification \cite{sannella-tarlecki}. The categorical semantics of database schemas as sketches is due to Johnson and Rosebrugh \cite{johnson-rosebrugh}; functorial data migration, with which this work is often conflated and from which it must be distinguished, is due to Spivak \cite{spivak2012,spivak2014}. The contribution here is at the level of transport and consequence:
\begin{enumerate}[leftmargin=1.4em,itemsep=0.2em]
  \item We transport the specification-gluing construction from algebraic specifications to AI-readable domain ontologies, fixing the formal object (finite-limit sketches), the bridge shape (spans through a shared sub-theory), and the integration (pushout and colimit of presentations).
  \item We prove that consistency equals commutativity (Proposition~\ref{prop:commute}) and that the model functor carries the theory-level colimit to the instance-level limit (Theorem~\ref{thm:soundness}), with an implementable governance rule as the strictness condition.
  \item We prove finishability (Proposition~\ref{prop:finish}): partial assemblies embed faithfully in the whole, which formalizes ``useful the moment it is written.''
  \item We diagnose the monolithic EDM and the OWL/RDF program, and identify LLM-legible sketches as what makes the construction operational.
\end{enumerate}
The intended venue is an industry or systems track. The paper is an opinionated architecture with a load-bearing formal core rather than a contribution to category theory.

\paragraph{What this paper is not.} The construction divides into a static half and a dynamic half, and this paper is deliberately the static one: composition that is correct at a fixed moment in organizational time. The dynamic half---continuous verification of the bridge annotations against live data, the bitemporal evolution of domains and bridges as regulation and acquisition move them, and the use of AI to keep the gluing current---is the subject of the companion practitioner volume \cite{semantic-enterprise}. We keep it out of scope here on purpose, because entangling it is exactly what would cost the soundness proof (Section~\ref{sec:soundness}) its elementary character. Where the static theory leaves a gap that the dynamic half closes, we say so and cite forward rather than gesture vaguely.

\section{Domains as finite-limit sketches}\label{sec:domains}

We need a formal object for ``a business domain's data model with its rules.'' We use the finite-limit sketch, which is exactly the categorical form of an entity--relationship--attribute model \cite{johnson-rosebrugh}. We introduce it in the reader's own terms.

\begin{definition}[Domain theory]\label{def:domain}
A \emph{domain} $D$ is a finite-limit sketch: a finite directed graph whose
\begin{itemize}[leftmargin=1.4em,itemsep=0.1em]
  \item \textbf{nodes} are the domain's \emph{entities} (sorts): $\Trade$, $\Client$, $\Wallet$, $\Inst$;
  \item \textbf{edges} are \emph{operations}---functional relationships read as maps---such as $\mathsf{owner}\colon \Wallet \to \Client$ (each wallet has one owning client) or $\mathsf{instr}\colon \Trade \to \Inst$ (each trade references one instrument); and
\end{itemize}
together with a finite set of \emph{imposed limit cones} expressing the domain's rules---among them the rules that make certain operations \emph{partial} or \emph{multivalued} (see below). A \emph{model} (instance) of $D$ is an assignment of a set to each entity and a function to each operation that sends every imposed cone to an actual limit cone in $\Set$. We write $\Mod(D)$ for the category of models of $D$ and their natural transformations (the maps of instances that commute with every operation).
\end{definition}

The theory is the schema-with-rules; a model is a concrete population of it, Tuesday's data. Three features of real domains motivate the sketch over plain equational theories, each one a place where a less expressive formalism would require the model to misdescribe the data.

\paragraph{Partiality.} A \emph{prospect} in a sales domain has no client record until it converts; a proprietary trading desk has no KYC \emph{customer} antecedent because it was never onboarded. The map ``prospect $\to$ client'' is therefore \emph{partial}, defined on a subset. A finite-limit sketch expresses this directly: the domain of definition is carved out as a limit (a monomorphism $\mathsf{Converted} \rightarrowtail \mathsf{Prospect}$ with a total map $\mathsf{Converted} \to \Client$), in place of a total operation that the data then violates. Partiality is pervasive in enterprise data, and the formalism represents it directly.

\paragraph{Multiplicity.} A client holds many wallets. The relationship $\Client$--$\Wallet$ is a function in the other direction, $\mathsf{owner}\colon \Wallet \to \Client$, or a span realizing a many-to-many tie. Sketches model both with the same limit machinery, without flattening the data.

\paragraph{Time and identity are ordinary attributes.} Two columns reliably disorient practitioners: the slowly-changing-dimension (SCD2) validity interval and the surrogate-key ``id fountain.'' In the static theory we model these as ordinary attributes---operations $\mathsf{validFrom},\mathsf{validTo}\colon \Trade \to \mathsf{Time}$ and $\mathsf{id}\colon \Trade \to \mathsf{Key}$---and give them no further status. The disorientation comes from giving them special logical status (temporal logics, identity criteria baked into the theory): these attributes feel load-bearing because they are the fields one reaches for when diagnosing a production incident or driving a replay, and that operational salience gets mistaken for ontological specialness. Bitemporality is a discipline applied across the layers---data, ontology, and the bridges between ontologies---and is developed in the companion practitioner volume \cite{semantic-enterprise}; the static composition results here stay clean because they do not entangle it. Treating valid-time and surrogate identity as attributes keeps the soundness proof (Section~\ref{sec:soundness}) elementary.

\begin{example}[A trading domain]\label{ex:trading}
The Trading domain has entities $\Trade,\Inst,\mathsf{Market},\mathsf{Desk}$ and operations $\mathsf{instr}\colon\Trade\to\Inst$, $\mathsf{listing}\colon\Inst\to\mathsf{Market}$, $\mathsf{desk}\colon\Trade\to\mathsf{Desk}$. Composition is path-following: $\mathsf{listing}\circ\mathsf{instr}\colon\Trade\to\mathsf{Market}$ gives the market of any trade, and the model must make this composite well defined (an imposed condition). A model populates these entities with the actual trades, instruments, markets and desks, and the actual reference functions among them.
\end{example}

Every well-modeled domain already carries this structure. Naming it leaves the domain unchanged and makes available the results about what happens when two domains meet.

\section{Bridges as spans}\label{sec:bridges}

Domains share concepts. A trade in Trading triggers a journal entry in Finance; a counterparty in Trading is a customer in Compliance; an instrument in Trading is a financial product in Finance. The central modeling decision of this paper concerns the \emph{shape} of the connection between two domains, and that shape is a span rather than a functor.

A functor $D_1 \to D_2$ would map all of $D_1$ into $D_2$, sending every Trading entity to a Finance entity, which makes Trading a part of Finance. Integration calls for the opposite: the two domains glued along the concepts they genuinely share, with everything else kept disjoint. The shape that does this is a \emph{span}.

\begin{definition}[Bridge]\label{def:bridge}
A \emph{bridge} between domains $D_1$ and $D_2$ is a span of sketch morphisms
\[
\begin{tikzcd}[column sep=large]
 & S \arrow[dl,"f_1"'] \arrow[dr,"f_2"] & \\
D_1 & & D_2
\end{tikzcd}
\]
where $S$ is the \emph{shared sub-theory}: the entities, operations, and rules the two domains hold in common. Each leg $f_i$ carries an \emph{annotation}: which constructs of $S$ are preserved into $D_i$, which entities of $D_i$ have no antecedent in $S$ (partiality of the integration), and the provenance of the mapping (who or what attested it, and when it was last verified).
\end{definition}

The shared sub-theory is governed by a rule that is simultaneously a design discipline and, as we will see, the precise hypothesis of the soundness theorem.

\begin{quote}\itshape
Sharing rule: if two domains genuinely share a concept, their rules for it agree, and $S$ contains that concept with those rules. If their rules disagree, they do not share the concept: split it into two distinct concepts, one per domain, and let the bridge document how they relate as a documented partiality rather than an identity.
\end{quote}

The rule has teeth. Consider $\Inst$ in Trading and $\FinProd$ in Finance. If both obey the same identity rule (same legal-entity key, same notion of when two are equal), they are the same concept; $S$ contains a single entity $\mathsf{Asset}$ with that rule, and both legs are inclusions of it. If Finance additionally partitions financial products by accounting treatment in a way Trading does not recognize, that partition is \emph{not} shared: it lives in $D_2$ above $S$, outside the bridge, and the bridge records that the accounting partition is lost on the Trading side. Under the sharing rule, the legs $f_i$ are \emph{faithful embeddings of $S$ that strictly preserve its imposed cones}---there is no reconciliation of conflicting axioms to perform, because conflicts are resolved by splitting before the bridge is drawn. We call such a span \emph{clean}. Cleanliness is the condition that makes Theorem~\ref{thm:soundness} hold on the nose rather than up to coherent isomorphism.

\begin{example}[A clean bridge with documented loss]\label{ex:bridge}
Compliance has $\Cust$ with operations $\mathsf{legalEntity}\colon\Cust\to\mathsf{LEI}$ and $\mathsf{kycClass}\colon\Cust\to\mathsf{KYCClass}$. Trading has $\Cpty$ with $\mathsf{legalEntity}\colon\Cpty\to\mathsf{LEI}$. The shared sub-theory $S$ contains the entity carrying the legal-entity identity and its $\mathsf{legalEntity}$ map---the concept on which the two domains genuinely agree. The legs include $S$ into both. The annotation records two facts. First, $\mathsf{kycClass}$ is \emph{lost} across the bridge: it lives in Compliance above $S$ and has no Trading antecedent, and forcing it into Trading would corrupt the trading model. Second, the integration is \emph{partial}: proprietary desks appear in Trading as counterparties with no Compliance customer antecedent, because they were never KYC-onboarded. Both facts are properties of the span, recorded once and inherited by every integration built from it.
\end{example}

\section{Integration as pushout, multi-domain as colimit}\label{sec:pushout}

Given a bridge, the integrated model is the pushout: the smallest gluing that identifies the shared sub-theory and keeps everything else disjoint.

\begin{definition}[Integrated model of two domains]\label{def:integration}
The integrated model of a bridge $D_1 \xleftarrow{f_1} S \xrightarrow{f_2} D_2$ is the pushout $D_1 +_S D_2$, taken in the category of \emph{sketch presentations}:
\[
\begin{tikzcd}[column sep=large,row sep=large]
S \arrow[r,"f_2"] \arrow[d,"f_1"'] & D_2 \arrow[d,"\iota_2"] \\
D_1 \arrow[r,"\iota_1"'] & D_1 +_S D_2 \arrow[ul, phantom, "\ulcorner"', very near start]
\end{tikzcd}
\]
Concretely: take the disjoint union of the entities, operations, and imposed cones of $D_1$ and $D_2$, and quotient by identifying the image of $S$ under the two legs. The result is the sketch presented by all of $D_1$'s and $D_2$'s declarations, with the shared declarations counted once.
\end{definition}

\paragraph{Where the pushout is taken.} A reader who knows that pushouts of categories along a subcategory are badly behaved---the pushout in $\Cat$ is generated by formal composites and need not be the naive gluing---will ask whether Definition~\ref{def:integration} is well-posed. It is, because the pushout is taken not in $\Cat$ but in the category of \emph{presentations} (sketches), where it is computed syntactically: disjoint union of generators and relations, modulo the shared sub-presentation. This is standard in the algebraic-specification tradition \cite{sannella-tarlecki,pwc2005} and is where the construction was always correct. The syntactic pushout could, in principle, present a theory whose \emph{models} are not the naive gluing of model populations. Theorem~\ref{thm:soundness} establishes that this does not occur.

\begin{definition}[Integrated enterprise model]\label{def:colimit}
A bridge diagram $\mathfrak{D}$ is a finite diagram of domains and bridges (a finite set of domain sketches together with a finite set of spans among them). The \emph{integrated enterprise model} is the colimit $\operatorname{colim}\mathfrak{D}$ of this diagram in the category of sketch presentations. When $\mathfrak{D}$ is tree-shaped the colimit is computed by iterated pushout---glue domains pairwise along their bridges, then glue the results along their shared sub-theories, and so on. A bridge diagram in general carries cycles (Section~\ref{sec:commute}), and there the colimit additionally coequalizes the parallel paths around each cycle; presentations are cocomplete, so $\operatorname{colim}\mathfrak{D}$ exists in either case. Each domain comes with a canonical leg $\iota_i\colon D_i \to \operatorname{colim}\mathfrak{D}$ into the integration.
\end{definition}

\begin{example}[The running integration]\label{ex:running}
Trading (Example~\ref{ex:trading}) and Compliance (Example~\ref{ex:bridge}) are glued along the shared legal-entity concept $S$. The pushout keeps $\Trade,\Inst,\mathsf{Market},\mathsf{Desk}$ from Trading and $\Cust,\mathsf{KYCClass}$ from Compliance, identifies the legal-entity concept once, and adds nothing else. By Theorem~\ref{thm:soundness}, an instance of the integration is a Trading population and a Compliance population that agree on legal-entity identity---precisely the join one wants for ``which counterparties are KYC-classified how,'' with proprietary desks correctly appearing as counterparties that lack a customer record.
\end{example}

The integrated enterprise model is, by the universal property of the colimit, the \emph{initial} gluing: any other model receiving compatible maps from all the domains factors through it uniquely. Operationally, there is exactly one canonical cross-domain view per diagram, and any hand-built cross-domain mart either equals it or is wrong relative to the bridges it claims to respect. This is the formal content of the practitioner slogan that bespoke cross-domain joins reinvent, badly, a structure that already exists canonically.

\section{Consistency is commutativity}\label{sec:commute}

The colimit always exists; presentations are cocomplete. The question that matters operationally is whether each domain embeds into the integration without two of its distinct concepts being silently forced together. This is where multiple bridges interact, and it is where defects live.

Write the colimit as a presentation whose generators are the disjoint union of all domains' entities and operations, and whose relations are the identifications imposed by the bridge legs. The legs generate an equivalence relation $\sim$ on generators: $x \sim y$ when some chain of bridges identifies them. The colimit identifies exactly the $\sim$-classes.

\begin{definition}[Commuting diagram, witness]
A bridge diagram $\mathfrak{D}$ \emph{commutes} if for every parallel pair of bridge-paths $p,q\colon D_i \rightrightarrows D_k$ in the span-graph and every concept $c$ carried along both paths, the transported images coincide: $p(c)$ and $q(c)$ are the same generator of $D_k$. A \emph{witness} to non-commutativity is a triple $(c,p,q)$: a concept $c$ and two parallel paths whose induced identifications, followed through, equate two generators that their home domain distinguishes.
\end{definition}

The witnessing situation is the cyclic one below: a shared concept reached from Trading both directly (via Compliance) and indirectly (via Finance), where the two routes disagree.
\[
\begin{tikzcd}[column sep=3.2em, row sep=2.0em]
& \mathsf{Trading} \arrow[dl, "p"', dash] \arrow[dr, "q", dash] & \\
\mathsf{Compliance} \arrow[rr, dash] & & \mathsf{Finance}
\end{tikzcd}
\qquad
\raisebox{2.2em}{$\displaystyle p(c)\overset{?}{\sim}q(c)$}
\]
If the counterparty identity $c$ transported along $p$ (direct Trading--Compliance bridge) and along $q$ (Trading--Finance--Compliance) disagree, the colimit is forced to identify two distinct Trading counterparties, collapsing the domain; the witness $(c,p,q)$ names the single concept and the two bridges to repair.

\begin{proposition}[Consistency $=$ commutativity]\label{prop:commute}
Let $\mathfrak{D}$ be a finite bridge diagram of clean spans. The canonical legs $\iota_i\colon D_i \to \operatorname{colim}\mathfrak{D}$ are jointly faithful---no two distinct concepts of any single domain are identified in the integration---if and only if $\mathfrak{D}$ commutes. Moreover, when commutativity fails, the failure is exhibited by a witness $(c,p,q)$ computable from the diagram.
\end{proposition}

\begin{proof}
Each $\iota_i$ is the composite $D_i \hookrightarrow \bigsqcup_j D_j \twoheadrightarrow \operatorname{colim}\mathfrak{D}$, where the quotient is by $\sim$. So $\iota_i$ identifies two generators $a,b$ of $D_i$ iff $a \sim b$, i.e.\ iff some chain of bridge identifications connects $a$ to $b$. Since each leg of a clean span is a faithful embedding of the shared sub-theory and adds no relations among the non-shared generators of its domain, any chain connecting $a$ to $b$ must leave $D_i$ through a shared concept and return through another, i.e.\ it is a pair of bridge-paths $p,q$ from $D_i$ to some $D_k$ and back, meeting at a shared concept $c$ with $p(c)$ and $q(c)$ landing on $a$ and $b$ respectively.

$(\Leftarrow)$ Suppose $\mathfrak{D}$ commutes. Then for every such $c$ and every parallel pair $p,q$, the identifications agree, so the only generators forced equal are those declared equal by a single genuine sharing. Restricted to the generators of a fixed $D_i$, $\sim$ is therefore the identity: no chain equates two distinct concepts of $D_i$. Hence each $\iota_i$ is injective on generators and, since clean legs preserve the imposed cones strictly, faithful.

$(\Rightarrow)$ Contrapositive. Suppose $\mathfrak{D}$ does not commute, witnessed by $(c,p,q)$ with $p(c)$ and $q(c)$ distinct in their home domain $D_i$ yet equated by the two routings. Then $p(c) \sim q(c)$ holds in the colimit while $p(c)\neq q(c)$ in $D_i$, so $\iota_i$ is not injective: the legs are not jointly faithful.

The witness is computable: enumerate the (finitely many) bridge-paths between each ordered pair of domains in the finite span-graph; for each parallel pair and each shared concept on both, check whether the induced identifications agree. Disagreement yields $(c,p,q)$ directly.
\end{proof}

\begin{remark}[The operational payoff]
A failed integration does not implicate ``the model.'' It implicates one concept and two bridges. The witness is the smallest possible repair surface: you re-examine a single span, not a global schema. This converts integration correctness from an opinion defended in a meeting into a check that returns a counterexample, in the same sense that a failing test returns a counterexample rather than a verdict.
\end{remark}

\begin{example}[A witness and its repair: counterparty identity under prime brokerage]\label{ex:witness}
Extend the running integration with a third domain. \textsf{Finance} carries a booking entity $\mathsf{Account}$ with an operation $\mathsf{holderLei}\colon\mathsf{Account}\to\mathsf{LEI}$ giving the registered account-holder's legal entity. Two bridges now reach Compliance from Trading and both purport to carry the concept $c$, ``the legal entity by which Compliance knows this counterparty.''
\begin{itemize}[leftmargin=1.4em,itemsep=0.1em]
  \item $p$, the \emph{direct} Trading--Compliance bridge of Example~\ref{ex:bridge}, transports $c$ along the shared legal-entity identity: a counterparty's Compliance LEI is its \emph{own} $\mathsf{legalEntity}$.
  \item $q$, the \emph{routed} bridge Trading--Finance--Compliance, identifies a counterparty with its Finance booking account (Trading--Finance shares booking-account identity), then reads off that account's $\mathsf{holderLei}$.
\end{itemize}
The prime-brokerage fact breaks commutativity: a single master account aggregates trades for two distinct legal entities, so its $\mathsf{holderLei}$ is the lead or broker entity, not either underlying one. The two transports of $c$ therefore land on different generators---$p(c)=\mathsf{legalEntity}$ and $q(c)=\mathsf{holderLei}\circ\mathsf{bookingAcct}$, two distinct operations $\Cpty\to\mathsf{LEI}$ that the domain genuinely distinguishes. The diagram, by declaring both routes to carry the \emph{same} concept, forces $p(c)\sim q(c)$, collapsing own-identity into account-holder-identity; in the populated model this identifies counterparties that prime brokerage only co-books. The witness is the triple $(c,p,q)$: it names one concept and exactly the two bridges to repair, not ``the model.''

The repair is the sharing rule of Section~\ref{sec:bridges}. The two routes were never sharing counterparty identity: $p$ carries legal-entity identity and $q$ carries booking-account-holder identity, and under prime brokerage---a many-to-one consolidation of legal entities onto accounts---these are different concepts. Split them. Keep $\mathsf{legalEntity}\colon\Cpty\to\mathsf{LEI}$ as the concept shared with Compliance, and carry the booking path as a separate operation $\mathsf{bookingAcct}\colon\Cpty\to\mathsf{Account}$ with its own $\mathsf{holderLei}$, annotated as the consolidation it is. With the concepts split, neither route claims the other's generator, the parallel pair commutes, and Proposition~\ref{prop:commute} certifies the legs jointly faithful: the integration distinguishes the two counterparties while recording, through Finance, that they bill through one account.
\end{example}

\begin{remark}[Cost and the hierarchical escape]\label{rem:cost}
The check is decidable but its cost is path enumeration, worst-case exponential in a densely connected span-graph. Two facts make this acceptable. First, real bridge graphs are sparse: most domain pairs share nothing and are connected only through hubs, so the number of parallel paths is small. Second, the construction is hierarchical: clusters of tightly-related domains form sub-colimits that are themselves treated as single objects and glued at a higher level (colimits of colimits), bounding path length. The construction is most useful in the regime of roughly five to thirty domains; beyond that, hierarchical composition is mandatory, or one accepts that distant domain pairs are never directly checked.
\end{remark}

\section{Soundness: gluing schemas glues instances}\label{sec:soundness}

Definition~\ref{def:integration} glues at the level of \emph{theory}. A database reviewer is entitled to demand soundness at the level of \emph{data}: that an instance of the glued schema be a pair of instances of the two schemas agreeing on what they share, with nothing spuriously added or lost. The following theorem supplies it.

\begin{theorem}[Soundness of gluing]\label{thm:soundness}
Let $D_1 \xleftarrow{f_1} S \xrightarrow{f_2} D_2$ be a clean span and $D_1 +_S D_2$ its pushout (Definition~\ref{def:integration}). Then the model functor sends the pushout of theories to the pullback of instance categories:
\[
\Mod(D_1 +_S D_2)\;\cong\;\Mod(D_1)\;\times_{\Mod(S)}\;\Mod(D_2),
\]
where the maps into $\Mod(S)$ are restriction along the legs. Explicitly, an instance of the integrated model is precisely a pair $(M_1,M_2)$ of domain instances satisfying $M_1\restrict{S} = M_2\restrict{S}$, and a morphism of integrated instances is precisely a pair of domain-instance morphisms agreeing on $S$.
\end{theorem}

\begin{proof}
We argue directly at the level of sketches and their set-valued models; the elementary argument yields the strict isomorphism and exposes exactly where cleanliness is used. Afterwards we record the structural reason.

A model $M$ of $D_1 +_S D_2$ assigns a set to each entity and a function to each operation of the pushout presentation, subject to its imposed cones. By Definition~\ref{def:integration}, the entities of the pushout are $\mathrm{Ent}(D_1)\sqcup\mathrm{Ent}(D_2)$ modulo identification of $f_1(\mathrm{Ent}(S))$ with $f_2(\mathrm{Ent}(S))$, and likewise for operations; the imposed cones are the union of those of $D_1$ and $D_2$, with the cones of $S$ (present in both via the clean legs) counted once.

Define $M_1 := M\circ\iota_1$ and $M_2 := M\circ\iota_2$, the restrictions of $M$ along the canonical legs. Each $M_i$ assigns sets and functions to the entities and operations of $D_i$ and, since $\iota_i$ carries each imposed cone of $D_i$ to an imposed cone of the pushout which $M$ sends to a limit, $M_i$ sends every imposed cone of $D_i$ to a limit. Hence $M_i \in \Mod(D_i)$. Because the shared generators are identified in the pushout, $M$ assigns a single set to each shared entity and a single function to each shared operation; therefore $M_1$ and $M_2$ assign the same data to everything in $S$, i.e.\ $M_1\restrict{S} = M_2\restrict{S}$. Thus $M \mapsto (M_1,M_2)$ lands in the pullback.

Conversely, let $(M_1,M_2)$ satisfy $M_1\restrict{S}=M_2\restrict{S}$. Define $M$ on the pushout by $M\circ\iota_i := M_i$; this is well defined precisely because the only generators shared between the two definitions are those of $S$, on which $M_1$ and $M_2$ agree by hypothesis. (Here cleanliness is essential: the legs are faithful embeddings of $S$ that introduce no further identifications, so there is no generator outside $S$ on which the two definitions could disagree, and no imposed cone of $S$ that the two could satisfy incompatibly.) The imposed cones of the pushout are those of $D_1$ and of $D_2$; $M$ sends the former to limits because $M_1$ does and the latter because $M_2$ does, and the shared cones of $S$ are satisfied because $M_1\restrict{S}=M_2\restrict{S}$ does. Hence $M\in\Mod(D_1+_S D_2)$.

These assignments are mutually inverse by construction. The same argument applied to natural transformations---a natural transformation out of the pushout is a family of components on each entity commuting with each operation, hence a pair of natural transformations agreeing on $S$---shows the correspondence is functorial and bijective on morphisms. Therefore $\Mod(D_1+_S D_2)\cong \Mod(D_1)\times_{\Mod(S)}\Mod(D_2)$ as categories.
\end{proof}

\begin{remark}[The structural reason, for the categorically inclined]
The elementary proof is an instance of a general fact. The functor $T$ sending a sketch to the finite-limit theory it generates \cite{barr-wells} is left adjoint to the underlying-sketch functor and so preserves colimits; thus $T(D_1+_SD_2)\cong T(D_1)+_{T(S)}T(D_2)$. The model functor $\Mod(-)=\Lex(T(-),\Set)$ is, in each argument, a hom-like functor into $\Set$ and so sends colimits of theories to limits of model categories. The clean-span hypothesis upgrades the resulting pseudo-pullback (instances agreeing on $S$ up to coherent isomorphism) to a strict pullback (agreeing on the nose). The elementary argument is given here because it is self-contained and because it locates the strictness in a governance rule a practitioner can enforce rather than in a technical assumption.
\end{remark}

\begin{corollary}[Multi-domain soundness]\label{cor:multi}
For a finite bridge diagram $\mathfrak{D}$ of clean spans,
\[
\Mod(\operatorname{colim}\mathfrak{D})\;\cong\;\lim \Mod(\mathfrak{D}),
\]
the limit of the diagram of instance categories and restriction functors. An instance of the integrated enterprise model is exactly a tuple of domain instances that agree on every shared sub-theory.
\end{corollary}

\begin{proof}
Apply the structural remark (after Theorem~\ref{thm:soundness}) to the whole diagram at once: the theory functor $T$ is a left adjoint, so it preserves the colimit, $T(\operatorname{colim}\mathfrak{D})\cong\operatorname{colim} T(\mathfrak{D})$; and $\Mod(-)=\Lex(T(-),\Set)$, hom-like into $\Set$ in each argument, sends that colimit of theories to the limit of model categories, with the clean-span hypothesis upgrading the pseudo-limit to the strict limit exactly as in Theorem~\ref{thm:soundness}. This argument is uniform in the shape of $\mathfrak{D}$ and in particular covers the cyclic diagrams of Section~\ref{sec:commute}, where the colimit coequalizes parallel paths and $\Mod$ correspondingly equalizes them. When $\mathfrak{D}$ is tree-shaped the conclusion is also visible elementarily: the colimit is built by iterated pushout (Definition~\ref{def:colimit}), $\Mod$ sends each pushout to a pullback by Theorem~\ref{thm:soundness}, and the pullbacks compose to the limit.
\end{proof}

A row in the integrated view is a tuple of domain rows that agree where the domains overlap, so the view adds no record that the domains do not jointly support. Where a domain holds data above a shared sub-theory (Finance's accounting partition, Compliance's KYC class), that data appears in the integration and is absent from the others, as the bridge annotations declared. The colimit adds no unwarranted identification and respects every genuine one; it is the canonical view determined by the diagram.

\section{Finishability: parts are useful before the whole exists}\label{sec:finish}

The monolithic EDM's defining vice is that it has no partial value: an incomplete global schema is a model of nothing. The colimit, presented as a diagram, has the opposite property, and this is the conceptual heart of the paper.

\begin{proposition}[Locality / finishability]\label{prop:finish}
Let $\mathfrak{E}\subseteq\mathfrak{D}$ be a sub-diagram (a subset of domains together with the bridges among them). There is a canonical comparison morphism $\operatorname{colim}\mathfrak{E}\to\operatorname{colim}\mathfrak{D}$ commuting with the legs. If $\mathfrak{D}$ commutes, this morphism is a faithful embedding. Dually, restriction $\Mod(\operatorname{colim}\mathfrak{D})\to\Mod(\operatorname{colim}\mathfrak{E})$ is the forgetful functor that sends a global instance to its restriction to the sub-assembly.
\end{proposition}

\begin{proof}
The inclusion of diagrams $\mathfrak{E}\hookrightarrow\mathfrak{D}$ induces, by the universal property of the colimit, a unique comparison map $\operatorname{colim}\mathfrak{E}\to\operatorname{colim}\mathfrak{D}$ commuting with all legs. Faithfulness under commutativity is Proposition~\ref{prop:commute} applied within $\mathfrak{D}$: no two distinct concepts appearing already in $\mathfrak{E}$ are identified in $\operatorname{colim}\mathfrak{D}$, so the comparison adds no identifications. The dual statement is Corollary~\ref{cor:multi}: $\Mod$ sends the comparison to a projection out of the limit, which is restriction.
\end{proof}

The consequences are the properties the monolith lacks.

\emph{A bridge is usable as soon as it is written.} A two-domain sub-assembly is a faithful part of the eventual whole, so an agent or analyst working across those two domains sees the two local ontologies and the explicit bridge between them, and the answers computed there are correct in the full integration. No wait for the other domains is required.

\emph{Partial work survives later additions.} Adding a domain or a bridge extends the colimit, and by Proposition~\ref{prop:finish} the earlier sub-assembly still embeds faithfully, so work already in production stays correct. The fear that motivates ``freeze the model before anyone uses it''---that later additions will retroactively break earlier work---does not hold under commutativity.

\emph{Gaps are visible and localized.} A cross-domain question whose two domains are not yet bridged fails by surfacing a missing bridge rather than by returning a wrong number. The absence is a typed hole the diagram can name, and naming it specifies the next piece of work.

Finishability here means correctness and usability at every stage of incompleteness, rather than eventual completeness. The colimit over a hundred domains may never be fully assembled; the colimit over the seven that matter this quarter is in production and correct.

\section{Why this, and not the monolith or OWL}\label{sec:owl}

We can now state the diagnosis formally.

\paragraph{The monolithic EDM.} To ``build the enterprise data model'' in the traditional sense is to author $\operatorname{colim}\mathfrak{D}$ directly: to write down by hand the single global schema the diagram would have produced. This has two structural defects. It violates locality, because a directly-authored global schema has no faithful sub-objects to ship early: Proposition~\ref{prop:finish} applies only to an object presented as a colimit of a diagram one possesses, and the monolith is not presented that way. And it misplaces correctness, because a directly-authored monolith is checked for global consistency, a property of the whole that on failure does not localize to a span or return a witness. The method poses an intractable question (``is the entire model consistent and complete?'') in place of a decidable one (``do these two bridges agree on this concept?'', Proposition~\ref{prop:commute}). The slowness of core domains sharpens the point: the monolith fails on stationary input, so motion cannot be the cause.

\paragraph{OWL and RDF.} The Semantic Web program is the most rigorous attempt to do better. Its open-world assumption refuses to demand global completeness, which is the locality of Proposition~\ref{prop:finish}. The implementation then reinstated the monolith: meaning was encoded in description logic over a merged global graph, and correctness became logical consistency of that graph, a global non-localizing check of the kind that defeats the monolithic EDM. The expressive ambition (full first-order-flavored axioms) made consistency checking expensive and authored ontologies brittle, and the formalism intended to make meaning machine-processable left it humanly unauthorable at enterprise scale, so the artifacts that shipped were the human-readable glosses that had dropped the formal content. The categorical reframing retains the open-world locality and drops the description-logic monolith: correctness is commutativity of a diagram, which is local and returns a witness, in place of consistency of a graph, which is global and opaque.

\paragraph{Why now: the objects are legible.} The construction is operational in 2026 and was not in 2006 for a non-mathematical reason. The objects to be glued, namely the domain theories and the annotations on the bridges, can now be authored and maintained as structured, prose-anchored sketches that a capable language model reads directly: entities, operations and cones as machine-checkable structure, conventions, edge cases and documented losses as prose the model reasons over. The expensive, brittle markup that sank the OWL program is replaced by an artifact that is cheap to write, legible to humans, and consumable by the machine that is now the primary reader. The mathematics was available two decades ago; its objects have only recently become writable.

Listing~\ref{lst:sketch} makes this concrete: the clean bridge of Example~\ref{ex:bridge} written as a prose-anchored sketch. The structured part---entities, operations, and the imposed cone on the shared concept---is machine-checkable: a tool validates that the legs preserve the cone and that the equivalence is generated only by declared sharing (Proposition~\ref{prop:commute}). The prose part---the documented loss, the partiality, the provenance and last-verified date---is what a capable language model reasons over directly, and is exactly the content that description logic forced into unauthorable markup or dropped to a human-readable gloss. The object that was prohibitive to write in OWL is a page a practitioner and a model both read.

\begin{lstlisting}[style=sketch,caption={The clean bridge of Example~\ref{ex:bridge} as a prose-anchored sketch: machine-checkable structure above, model-legible prose below.},label={lst:sketch}]
bridge trading.Counterparty <-> compliance.Customer

shared S:
  entity Party
    op legalEntity : Party -> LEI
    cone  identity-on(legalEntity)        # LEI is the shared identity key

legs:
  into Trading:    Party |-> Counterparty,  legalEntity |-> cpty.legalEntity
  into Compliance: Party |-> Customer,       legalEntity |-> cust.legalEntity

annotations:
  lost:    Customer.kycClass               # no Trading antecedent; forcing it
                                           #   through would corrupt the trading model
  partial: proprietary desks appear in Trading as Counterparty with no
           Customer antecedent             # internal entities, never KYC-onboarded
  provenance:
    attested:      {trading-agent, compliance-agent}
    last_verified: 2026-05-31              # re-checked against live data each scan
\end{lstlisting}

\section{Related work}\label{sec:related}

\textbf{The parent.} The pushout-union of specifications, and of the proofs that inhabit them, within a constructive type-theoretic setting is the author's prior work \cite{cpw2000,pwc2005}: a calculus whose structural rules are renaming, hiding, and union, the last being the pushout. The present paper transports that construction from specifications of programs to the integration of domain data models, where the gluing is cheaper to state and the operational reading (consistency as commutativity, finishability as faithful sub-colimits) is the new content.

\textbf{The logical home.} Institutions \cite{goguen-burstall} provide the general theory of logical systems glued by truth-preserving maps; the legs of our spans are institution morphisms made into first-class governance artifacts carrying their annotations. Structured algebraic specification \cite{sannella-tarlecki} is the source of gluing-as-pushout at the level of presentations, which is exactly where we take the pushout (Section~\ref{sec:pushout}).

\textbf{The database-shaped precedent.} Johnson and Rosebrugh \cite{johnson-rosebrugh} give the categorical semantics of entity--relationship--attribute models as sketches with set-valued models; our Definition~\ref{def:domain} is their correspondence, used as the formal object for a business domain.

\textbf{The cousin.} Functorial data migration \cite{spivak2012,spivak2014} models schemas as categories and instances as functors, and moves instances along a fixed schema functor via the migration adjunction. It is frequently conflated with the present construction and differs from it: the migration setting moves instances along a given map, whereas this paper composes and integrates the schemas themselves, by spans and colimits, with a soundness result on instances. Spivak's work is the instance-level cousin; the specification-composition tradition is the schema-level parent.

\textbf{Evolution and matching.} Double-pushout rewriting \cite{ehrig} provides the calculus for the cases in which domains do change (Section~\ref{sec:evolution}). The problem of keeping cross-domain mappings consistent as ontologies evolve is the subject of the ontology-matching and mapping-maintenance literature \cite{euzenat-shvaiko}; our commutativity check is a categorical criterion for the consistency such work seeks to maintain.

\textbf{The practitioner foils.} Dimensional modeling \cite{kimball} and the corporate information factory contribute durable principles (grain, conformed dimensions, historicization) while assuming a total schema; data mesh \cite{dehghani} correctly distributes domain ownership but underspecifies the composition mechanism. The colimit-of-domains is the missing composition layer: it distributes ownership (each domain is its own sketch) while supplying the canonical, sound integration the mesh leaves to ad-hoc joins.

\section{Evolution, briefly}\label{sec:evolution}

Schema evolution is often treated as the central difficulty of enterprise integration. Because core business domains change slowly, the dynamic case is not where the value lies, and foregrounding it concedes the moving-target excuse this paper rejects. We therefore record the evolution calculus as a remark rather than a pillar.

A schema change is a double-pushout rewrite rule $L \xleftarrow{} K \xrightarrow{} R$ with a match into the current model: the pushout-complement retires $L\setminus K$ and a pushout introduces $R\setminus K$ \cite{ehrig}. The gluing conditions become governance preconditions: the \emph{dangling condition} forbids retiring an entity while a live bridge still references it, and the \emph{identification condition} forbids merging two entities a downstream bridge distinguishes. A dated sequence of such rewrites is the model's history. The emphasis runs against agility. The rules guarantee that a disciplined rewrite cannot silently introduce a contradiction, and core domains rarely need rewriting, so the live concern is correctness and coverage on a target that holds still for long stretches, rather than pace of change. The failure this paper diagnoses is stasis with incompleteness.

\section{Threats to validity}\label{sec:threats}

\textbf{Pushouts of presentations versus of categories.} Our pushout is syntactic (Section~\ref{sec:pushout}); the naive intuition of ``gluing the categories'' is not what we compute, and would misbehave. Theorem~\ref{thm:soundness} is what licenses the syntactic gluing by showing it is sound on instances. We do not claim the category-level pushout; we claim the presentation-level one and its instance-level soundness.

\textbf{Strictness depends on the sharing rule.} Theorem~\ref{thm:soundness} gives a strict pullback under the clean-span hypothesis; without it one obtains a pseudo-pullback (agreement up to coherent isomorphism). The clean-span hypothesis is the enforceable governance rule of Section~\ref{sec:bridges}, but it is a hypothesis, and an organization that declares two genuinely-conflicting concepts ``the same'' violates it and loses strictness. The rule is sound discipline, not a theorem about human behavior.

\textbf{Annotations are asserted, not verified.} Nothing in the mathematics forces a bridge's declared losses and partialities to be true. The construction guarantees that \emph{if} the annotation is accurate, the integration is sound; it does not verify accuracy. Closing this gap requires tying each bridge to executable conformance checks on instances---continuous verification of the morphisms against live data---which is operational infrastructure outside the scope of the static theory, and the subject of the companion practitioner volume \cite{semantic-enterprise}.

\textbf{Cost in the dense regime.} The commutativity check is decidable, with cost in path enumeration (Remark~\ref{rem:cost}); beyond roughly thirty densely-connected domains, hierarchical composition is required or distant pairs go directly unchecked. We record this as a limitation of the approach.

\section{Conclusion}

The unified enterprise data model has been pursued for forty years and abandoned in nearly every instance, usually with the explanation that the enterprise moves too fast. Core domains move slowly, so that explanation does not hold, and the model failed for a structural reason: it was authored as a monolith and checked globally, where it should have been assembled as a colimit of a diagram and checked by the commutativity of its bridges. Named correctly, the construction is elementary, the specification-gluing of \cite{cpw2000,pwc2005} transported to a new domain, and it carries the properties the monolith lacked: a decidable consistency check that returns a witness (Proposition~\ref{prop:commute}), soundness from schemas to instances (Theorem~\ref{thm:soundness}), and finishability, under which every part is correct and usable before the whole exists (Proposition~\ref{prop:finish}). The Semantic Web program identified the open-world assumption and then encoded it as a description-logic monolith; that monolith is what this paper removes, and the objects it relies on instead, domain ontologies legible to both humans and language models, are now cheap to build.

The continuity with \cite{cpw2000} is worth stating directly, because it locates what is and is not new. That paper glued formal specifications along a shared sub-specification by the same pushout, with the structural operations of renaming, hiding, and union; the present paper applies that construction unchanged, replacing the specifications with domain ontologies and the set of models of a specification with the set of populations of a domain. The mechanism is identical: a span through what two objects share, a pushout that identifies the shared part and keeps the rest disjoint, a colimit that composes many such gluings. What differs is the domain and what the domain makes salient. In the 2000 setting the gluing served program synthesis, and composition was the object of study. Here the gluing serves data integration, the static composition is comparatively easy to state, and the load shifts onto two readings that the earlier setting did not need: that the gluing is sound at the level of populations (Theorem~\ref{thm:soundness}), and that partial assemblies are usable before the whole exists (Proposition~\ref{prop:finish}). The construction travels intact across twenty-five years and two domains; the consequences that matter are the ones the new domain forces into view. The mathematics was settled decades ago; its objects have only recently become writable.

\begin{thebibliography}{99}
\small
\bibitem{cpw2000} J.~N.~Crossley, I.~H.~Poernomo, and M.~Wirsing. Extraction of structured programs from specification proofs. In \emph{Recent Trends in Algebraic Development Techniques (WADT'99)}, LNCS 1827, pp.\ 419--437. Springer, 2000.
\bibitem{pwc2005} I.~H.~Poernomo, M.~Wirsing, and J.~N.~Crossley. \emph{Adapting Proofs-as-Programs: The Curry--Howard Protocol}. Monographs in Computer Science. Springer, 2005.
\bibitem{goguen-burstall} J.~A.~Goguen and R.~M.~Burstall. Institutions: abstract model theory for specification and programming. \emph{Journal of the ACM}, 39(1):95--146, 1992.
\bibitem{sannella-tarlecki} D.~Sannella and A.~Tarlecki. \emph{Foundations of Algebraic Specification and Formal Software Development}. Monographs in Theoretical Computer Science. Springer, 2012.
\bibitem{johnson-rosebrugh} M.~Johnson and R.~Rosebrugh. Entity-relationship-attribute designs and sketches. \emph{Theory and Applications of Categories}, 10(3):94--112, 2002.
\bibitem{spivak2012} D.~I.~Spivak. Functorial data migration. \emph{Information and Computation}, 217:31--51, 2012.
\bibitem{spivak2014} D.~I.~Spivak. \emph{Category Theory for the Sciences}. MIT Press, 2014.
\bibitem{ehrig} H.~Ehrig, K.~Ehrig, U.~Prange, and G.~Taentzer. \emph{Fundamentals of Algebraic Graph Transformation}. Springer, 2006.
\bibitem{euzenat-shvaiko} J.~Euzenat and P.~Shvaiko. \emph{Ontology Matching}, 2nd ed. Springer, 2013.
\bibitem{kimball} R.~Kimball and M.~Ross. \emph{The Data Warehouse Toolkit}, 3rd ed. Wiley, 2013.
\bibitem{dehghani} Z.~Dehghani. \emph{Data Mesh: Delivering Data-Driven Value at Scale}. O'Reilly, 2022.
\bibitem{semantic-enterprise} I.~Poernomo. \emph{The Semantic Enterprise: Data Architecture for the Agentic Era}. O'Reilly, forthcoming.
\bibitem{barr-wells} M.~Barr and C.~Wells. \emph{Toposes, Triples and Theories}. Grundlehren der mathematischen Wissenschaften 278. Springer, 1985.
\end{thebibliography}

\end{document}
